\subsection{序列层审计的周期核、同步核与异常核三分解}\label{subsec:conclusion-sequence-layer-period-sync-anomaly-audit-decomposition}
沿用小节 \ref{subsec:conclusion-periodic-fiber-cnf-complexity}、小节 \ref{subsec:conclusion-order-sync-torsion-prony-s5-defectline}、小节 \ref{subsec:conclusion-pressure-curve-s3-jordan-cumulant-universality} 与小节 \ref{subsec:conclusion-window6-mod2-k0-sync-replay-boundary-face} 的记号。前述结论链已经分别固定了四个彼此分散的接口：其一，cover 周期点、内禀 \(\zeta\) 与 determinant 层对歧义壳的常数级压缩与谱盲性；其二，未同步前缀尾部、扫描预算与 Hausdorff 余维之间的精确互反；其三，折叠多重度的最小线性实现维数由 follower ambiguity classes 唯一决定；其四，window-\(6\) 中局部精确反演与 anomaly 账本保真分属不同的二进资源尺度。把这四条链条并读后，可把固定分辨率的 sequence-layer audit 压缩为一条新的三分解原则：周期核、同步核与异常核彼此正交，且任一层都不能由其余两层吸收。

\begin{theorem}[序列层审计的周期核与同步核严格正交]\label{thm:conclusion-sequence-layer-periodic-sync-kernel-orthogonality}
设 \(m\ge 3\)。记 \(A_m\) 为 \(Y_m\) 的一个 follower-separated 右解析 cover 的邻接矩阵，并记相应右确定化表示为
\[
A_m^{\det}
=
\begin{pmatrix}
N_m & B_m\\
0 & A_m^{\mathrm{ess}}
\end{pmatrix},
\]
其中 \(N_m\) 对应歧义壳，\(A_m^{\mathrm{ess}}\) 对应单点骨架。再定义长度-\(n\) 周期压缩因子
\[
R_m(n):=
\frac{\Tr(A_m^n)}{\#\mathrm{Fix}(\sigma^n|_{Y_m})}
\qquad (n\ge 1).
\]
则成立：
\begin{enumerate}
  \item
  \[
  N_m^m=0,
  \qquad
  \det(I-zA_m^{\det})=\det(I-zA_m^{\mathrm{ess}});
  \]
  \item 对所有满足 \(\#\mathrm{Fix}(\sigma^n|_{Y_m})>0\) 的 \(n\)，
  \[
  1\le R_m(n)\le |\mathcal S_m|;
  \]
  \item 因而周期/Fredholm 数据在指数尺度上只感知单点骨架，而 cover 到 sofic 的损失至多是常数级；
  \item 若歧义壳非空，并且其最大谱半径强连通分量周期为 \(p\)，则对 \(Y_m\) 的最大熵测度 \(\nu_m\) 与同步时间 \(T_m\) 有
  \[
  \nu_m(T_m>n)=c_{n\bmod p}(m)e^{-\varepsilon_m n}(1+o(1)),
  \qquad
  \varepsilon_m=\log 2-\log \rho_{\mathrm{amb}}(m),
  \]
  其中 \(\rho_{\mathrm{amb}}(m)\) 为该歧义壳主块的 Perron--Frobenius 谱半径。
\end{enumerate}
因此，周期核与同步核在结构上严格正交：前者由 \(A_m^{\mathrm{ess}}\) 决定，后者由歧义壳的 Perron--Frobenius 数据决定。
\end{theorem}
\begin{proof}
由定理 \ref{thm:conclusion-ambiguity-shell-spectral-invisibility}，歧义壳块满足
\[
N_m^m=0,
\qquad
\det(I-zA_m^{\det})=\det(I-zA_m^{\mathrm{ess}}).
\]
因此，确定化表示的 determinant、\(\zeta\)-行列式与一切只锚定非零谱的 Fredholm 型不变量都只依赖于单点骨架。

另一方面，命题 \ref{prop:Ym-zeta-leading-pole} 与命题 \ref{prop:periodic-count-sandwich} 已给出：cover 周期点计数 \(\Tr(A_m^n)\) 与 sofic 因子周期点计数 \(\#\mathrm{Fix}(\sigma^n|_{Y_m})\) 具有同一主极点模长，且其差异至多为常数因子。推论 \ref{cor:Ym-periodic-compression-ratio} 则把该常数级缺陷精确写为
\[
1\le R_m(n)\le |\mathcal S_m|.
\]
故周期/Fredholm 通道在指数尺度上只看到单点骨架，而看不到歧义壳。

同步层则不同。由定理 \ref{thm:Ym-ambiguity-spectrum-law}，只要歧义壳主谱块非空，其同步尾便满足
\[
\nu_m(T_m>n)=c_{n\bmod p}(m)e^{-\varepsilon_m n}(1+o(1)),
\qquad
\varepsilon_m=\log 2-\log \rho_{\mathrm{amb}}(m).
\]
于是同步预算完全由歧义壳的 Perron--Frobenius 数据与周期调制控制。两部分合并，即得周期核与同步核的严格正交分裂。证毕。
\end{proof}

\begin{theorem}[清晰度成本恰等于歧义壳余维的倒数尺度]\label{thm:conclusion-sequence-layer-clarity-cost-codim-reciprocal}
设
\[
Y_m^{\mathrm{amb}}:=\mathcal Y_{m,\ge 2}
\]
为“始终至少保留两重候选”的歧义壳，并假定其最大谱半径强连通分量唯一且 primitive。记
\[
\varepsilon_m:=\varepsilon_{m,\ge 2},
\qquad
c_m:=c_{m,\ge 2},
\qquad
n_\eta:=n_{m,\eta}^{(2)}.
\]
则当 \(\eta\downarrow 0\) 时，
\[
n_\eta
=
\frac{1}{\varepsilon_m}\log\frac{c_m}{\eta}+O(1)
=
\frac{1}{(\log 2)\,\operatorname{codim}(Y_m^{\mathrm{amb}})}
\log\frac{c_m}{\eta}+O(1).
\]
因此，清晰度预算的主导系数并不是外加经验常数，而恰是歧义壳前缀超度量余维的严格倒数。
\end{theorem}
\begin{proof}
由定理 \ref{thm:conclusion-scan-budget-codim-duality}，在 \(k=2\) 的阈值层上有
\[
n_\eta
=
\frac{1}{\varepsilon_m}\log\frac{1}{\eta}
+
\frac{1}{\varepsilon_m}\log c_m
+
O(1)
=
\frac{1}{\varepsilon_m}\log\frac{c_m}{\eta}+O(1).
\]
再由定理 \ref{thm:conclusion-threshold-ambiguity-hausdorff-codim-ladder}，
\[
\operatorname{codim}(Y_m^{\mathrm{amb}})
=
\frac{\varepsilon_m}{\log 2}.
\]
代回前式即得
\[
n_\eta
=
\frac{1}{(\log 2)\,\operatorname{codim}(Y_m^{\mathrm{amb}})}
\log\frac{c_m}{\eta}+O(1).
\]
证毕。
\end{proof}

\begin{theorem}[常数级周期缺陷不能替代对数级清晰度预算]\label{thm:conclusion-sequence-layer-constant-periodic-defect-cannot-replace-clarity-budget}
\leanverified{paper\_conclusion\_sequence\_layer\_constant\_periodic\_defect\_cannot\_replace\_clarity\_budget}
沿用上面的记号，并定义常数级 sofic 缺陷
\[
\delta_m:=\limsup_{n\to\infty}\log R_m(n).
\]
若 \(Y_m^{\mathrm{amb}}\neq\varnothing\)，则
\[
0\le \delta_m\le \log|\mathcal S_m|
\]
且
\[
\lim_{\eta\downarrow 0}\frac{\delta_m}{n_\eta}=0.
\]
因此，周期层所记录的压缩损失始终只是常数级，而清晰度恢复所需的同步预算在 \(\eta\downarrow 0\) 时必然呈对数发散；任何仅依赖周期/Fredholm 缺陷的审计都不可能替代同步预算。
\end{theorem}
\begin{proof}
由推论 \ref{cor:Ym-periodic-compression-ratio}，
\[
0\le \delta_m\le \log|\mathcal S_m|<\infty.
\]
另一方面，由定理 \ref{thm:conclusion-sequence-layer-clarity-cost-codim-reciprocal}，
\[
n_\eta=\frac{1}{\varepsilon_m}\log\frac{c_m}{\eta}+O(1)
\xrightarrow[\eta\downarrow 0]{}\infty.
\]
故有限常数 \(\delta_m\) 除以 \(n_\eta\) 的极限必为 \(0\)。证毕。
\end{proof}

\begin{theorem}[精确纤维计数的最小状态维数由右歧义类决定，且不能由周期层单独恢复]\label{thm:conclusion-sequence-layer-fiber-counting-minimal-state-not-periodically-recoverable}
设折叠系统存在有限窗逆码，并可由 follower-separated 右解析表示逐位更新全部合法预像数。则存在有限状态集 \(S\)、边界向量 \(u,v\in\NN^S\) 与非负整数矩阵 \(T_0,T_1\) 使
\[
d_m(x)=u^\top T_{x_1}T_{x_2}\cdots T_{x_m}v
\qquad
(x=x_1\cdots x_m\in X_m).
\]
并且在全部此类线性实现中，最小可能维数恰等于 follower ambiguity classes 的基数。特别地，该极小维数是由右歧义结构唯一决定的规范量，而不能由内禀周期/\(\zeta\)/Fredholm 层单独恢复。
\end{theorem}
\begin{proof}
矩阵积表示
\[
d_m(x)=u^\top T_{x_1}\cdots T_{x_m}v
\]
以及“最小维数等于 follower ambiguity classes 的基数”正是定理 \ref{thm:conclusion-fold-inversecode-hilbert-recognizable} 的内容。

另一方面，由定理 \ref{thm:conclusion-sequence-layer-periodic-sync-kernel-orthogonality}，周期/\(\zeta\)/Fredholm 通道只感知单点骨架，而对歧义壳的同步结构失明；但 follower ambiguity classes 恰记录“给定输出前缀之后仍可能的候选集合如何分裂”的最小状态记忆，属于歧义层而非周期骨架。故该极小维数不能由周期层单独恢复。证毕。
\end{proof}

\begin{corollary}[window-\texorpdfstring{$6$}{6} 的十九比特正交裂分]\label{cor:conclusion-sequence-layer-window6-nineteen-bit-orthogonal-splitting}
在 window-\(6\) 的 bin-fold 口径下，局部精确反演所需最小位数与异常完备读出位数满足
\[
b_{\mathrm{inv}}=2,
\qquad
b_{\mathrm{anom}}=21,
\qquad
b_{\mathrm{anom}}-b_{\mathrm{inv}}=19,
\qquad
\frac{2^{b_{\mathrm{anom}}}}{2^{b_{\mathrm{inv}}}}=2^{19}.
\]
\leanverified{paper\_conclusion\_sequence\_layer\_window6\_nineteen\_bit\_orthogonal\_splitting}
因此，window-\(6\) 中至少存在 \(19\) 个独立异常比特，无法压缩进任何局部精确反演标签。
\end{corollary}
\begin{proof}
由定理 \ref{thm:conclusion-window6-inversion-anomaly-information-separation}，
\[
b_{\mathrm{inv}}=2,
\qquad
b_{\mathrm{anom}}=21.
\]
相减即得
\[
b_{\mathrm{anom}}-b_{\mathrm{inv}}=19,
\]
而相应寄存器容量之比为
\[
\frac{2^{b_{\mathrm{anom}}}}{2^{b_{\mathrm{inv}}}}
=
2^{21-2}
=
2^{19}.
\]
证毕。
\end{proof}

\begin{conjecture}[固定分辨率下 sequence-layer audit 的最小完备签名为三元组]\label{conj:conclusion-sequence-layer-minimal-complete-audit-signature-triple}
对每个固定分辨率 \(m\)，定义
\[
\mathfrak S_m
:=
\bigl(
\zeta_{Y_m}^{\mathrm{intr}},
\ U_m(z),
\ (b_{\mathrm{inv}}(m),b_{\mathrm{anom}}(m))
\bigr),
\]
其中
\[
U_m(z):=\sum_{n\ge 0}N_m^{\mathrm{unsync}}(n)z^n.
\]
则 \(\mathfrak S_m\) 应构成 fixed-resolution sequence-layer audit 的最小完备签名；亦即，任何删去其中一项的二元统计都不再普适完备。其分量分别负责：
\begin{enumerate}
  \item \(\zeta_{Y_m}^{\mathrm{intr}}\) 负责全部内禀周期/Fredholm 信息；
  \item \(U_m(z)\) 负责歧义壳的主奇点星簇、周期调制与清晰度预算；
  \item \((b_{\mathrm{inv}}(m),b_{\mathrm{anom}}(m))\) 负责局部精确反演与异常账本之间不可压缩的有限尺度裂分。
\end{enumerate}
其根据是三层刚性的：
\begin{enumerate}
  \item 定理 \ref{thm:conclusion-sequence-layer-periodic-sync-kernel-orthogonality} 表明 \(\zeta_{Y_m}^{\mathrm{intr}}\) 与同步尾机制严格正交；
  \item 定理 \ref{thm:conclusion-sequence-layer-clarity-cost-codim-reciprocal} 与定理 \ref{thm:conclusion-ambiguity-shell-periodic-amplitude-obstruction} 表明 \(U_m(z)\) 不仅承载扫描预算，而且承载歧义壳主奇点的周期振幅与星簇几何；
  \item 定理 \ref{thm:conclusion-sequence-layer-fiber-counting-minimal-state-not-periodically-recoverable} 表明右歧义类的极小状态维数不能由周期层吸收；
  \item 推论 \ref{cor:conclusion-sequence-layer-window6-nineteen-bit-orthogonal-splitting} 表明，即便局部精确反演已知，异常层仍至少留下一个独立的十九比特裂分。
\end{enumerate}
若该猜想成立，则固定分辨率下的 sequence-layer audit 不再只是若干充分统计量的并置，而会获得一个真正尖锐的最小签名。
\end{conjecture}

综上，序列层审计在固定分辨率上可压缩为
\[
\text{周期可见层}
\ \oplus\
\text{同步可见层}
\ \oplus\
\text{异常可见层},
\]
并且这三层并非“同一损失的不同坐标”，而是三个服从不同组合律、不同量纲与不同极小性原则的正交通道：周期层只承担骨架谱与 Fredholm 证书；同步层承担歧义壳的指数尾、余维与主奇点星簇；异常层则承担局部精确反演之外仍然存活的有限二进账本。因而，对本文的 fixed-resolution sequence-layer audit 而言，任何试图把这三层重新压缩到单一标量或单一有限寄存器中的方案，都会在同步预算或 anomaly 账本处发生不可修复的失真。

\endinput
